\documentclass[12pt,a4paper]{article}
\usepackage[utf8]{inputenc}
\usepackage[spanish]{babel}
\usepackage{amsmath}
\usepackage{amsfonts}
\usepackage{amssymb}
\usepackage{geometry}
\usepackage{graphicx}
\usepackage{titlesec}
\usepackage{parskip}
\usepackage{fancyhdr}
\usepackage{enumitem}

\geometry{margin=1in}

% Configuración de encabezado y pie de página
\pagestyle{fancy}
\fancyhf{}
\fancyhead[L]{\textbf{Examen de Lectura}}
\fancyhead[C]{\textbf{Control 1}}
\fancyhead[R]{\today}
\fancyfoot[C]{\thepage}

% Formato de títulos
\titleformat{\section}
{\normalfont\Large\bfseries}{\thesection}{1em}{}
\titleformat{\subsection}
{\normalfont\large\bfseries}{\thesubsection}{1em}{}

% Título principal
\title{\vspace{-1cm}
\huge\textbf{Control de Lectura N°1} \\
\Large\textbf{The Rise of Artificial Intelligence in Healthcare Applications} \\
\vspace{0.5cm}
\large{Bohr \& Memarzadeh (2020)}}
\author{\textbf{Armando Arredondo Valle} \\
\textit{Magíster en Inteligencia Artificial} \\
\textit{Pontificia Universidad Católica de Chile}}
\date{01 de abril de 2026}

\begin{document}

\maketitle
\thispagestyle{fancy}

\vspace{0.5cm}

\vspace{0.5cm}
\hrule
\vspace{0.5cm}

\section*{Pregunta N°1}
\textbf{¿Cómo explica la lectura que la IAaaa está ayudando a transformar la salud desde un modelo reactivo hacia uno más preventivo y personalizado?}

\subsection*{Respuesta}
A lo largo de la lectura se plantea justamente que la Inteligencia Artificial se está forjando como un paradigma fundamental en lo que respecta a la atención sanitaria, pudiendo incluso ser así que desplace el enfoque donde lo que podemos llamar como un "modelo reactivo" que en este caso va centrado en el tratamiento de enfermedades de igual forma al diagnóstico de las mismas. Ayudando así incluso que la transición que se maneje por parte de la IA permita que se realice de forma correcta.\\

Pudiendo ser así incluso que se permita a través de dispositivos que utilicen IA relacionados al monitoreo continuo nos permitirá identificar desviaciones que pueden ocurrir en parámetros en lo que respecta a la normalidad, ayudando incluso a tener intervenciones lo más pronto posibles, permitiendo tener reducciones en lo que podría ser incluso hospitalizaciones.


De igual forma nos habla sobre como la medicina de precisión se maneja por este nuevo modelo, siendo así que la IA permite analizar los datos individuales, tales como variaciones genómicas o incluso perfiles inmunológicos, donde el uso de la biología individual va a ser clave en todas las etapas del viaje médico.


Así mismo, nos permite hablar del empoderamiento del paciente, dando incluso énfasis en el aqcceso a registros propios de saludos, tales como los registros personales inteligentes así como la conexión de aplicaciones de monitoreo, donde cada uno mantiene el control de forma activa con la finalidad de tener este bienestar.


Siendo así que habla que puede llegar a tener implicancias económicas significativas, incluso con la posibilidad de reducir los costos anuales de salud, o de igual forma prevenir enfermedades.


\vspace{0.5cm}
\hrule
\vspace{0.5cm}

\section*{Pregunta N°2}
\textbf{Según la lectura, ¿qué aplicaciones de la IA en salud parecen hoy más prometedoras, y por qué se plantea que debe complementar el trabajo del personal clínico en vez de reemplazarlo?}

\subsection*{Respuesta}

A lo largo de la lectura se nota que la IA está abriendo puertas increíbles en salud, especialmente en lo que es la medicina de precisión. Es fascinante cómo empresas como Deep Genomics ya están cruzando datos genéticos para entender mejor las enfermedades, o cómo el deep learning está acelerando a pasos agigantados la creación de nuevos medicamentos, algo que antes tomaba años. También destaca mucho el tema del diagnóstico por imágenes; hoy los algoritmos pueden analizar radiografías o lesiones en la piel con una precisión que iguala o incluso supera a la de un ojo humano, llegando a niveles de acierto de casi el 93%.

Ahora, sobre por qué no va a reemplazar a los médicos, la razón es muy humana. La IA es una herramienta espectacular para procesar datos y automatizar lo administrativo, pero hay cosas que simplemente no puede copiar, como la empatía, la persuasión y el juicio clínico complejo. La idea no es que la máquina tome el control, sino que el personal clínico la use para liberarse de tareas repetitivas y así enfocarse en lo que realmente importa: el paciente. Al final, el riesgo no es que la IA reemplace al médico, sino que los médicos que no aprendan a usarla se queden atrás frente a quienes sí la aprovechen para mejorar su trabajo.

\vspace{0.5cm}
\hrule
\vspace{0.5cm}

\section*{Pregunta N°3}
\textbf{Pensando más allá de la salud, ¿cómo creen que la IA podría utilizarse en la prevención del delito? Mencionen una posible aplicación y un riesgo ético asociado.}

\subsection*{Respuesta}
Así como en salud se usan datos para predecir si un paciente volverá a enfermarse, en seguridad se podrían usar algoritmos para analizar patrones históricos, zonas geográficas y horarios para identificar dónde es más probable que ocurra un delito. Ciudades como Los Ángeles ya han probado estos sistemas para desplegar patrullas de forma más inteligente, tratando de actuar antes de que las cosas pasen en lugar de solo reaccionar al delito.

Sin embargo, el gran problema aquí es el sesgo de los datos. Si el sistema se entrena con información que ya viene de una policía que, por ejemplo, vigila más a ciertos grupos étnicos o barrios pobres, la IA solo va a repetir y aumentar esa discriminación. Se crea un círculo vicioso: el algoritmo dice que un barrio es "peligroso", la policía va más para allá, hace más detenciones, y eso confirma un sesgo del algoritmo.

\vspace{0.5cm}
\hrule
\vspace{0.5cm}

\section*{Pregunta N°4}
\textbf{A partir de estas dos primeras clases, ¿cuáles sienten que han sido los aprendizajes o ideas más importantes que se llevan hasta ahora?}

\subsection*{Respuesta}

Mayormente puedo decir que el manejo de la información de forma correcta. Es importante definir que la información a fin de cuentas tiene un origen y un fin, no podemos simplemente dejarla pasar así como así es necesario que como desarrolladores tengamos que tener cursos de ética, y que a fin de cuentas la información que planteemos y el uso que le demos tiene que ser importante y claramente tenga un impacto.

\end{document}